\section{Related Works}

% \textcolor{blue}{
% Potential logic for related works section
% \begin{itemize}
%     \item World Models for Planning and RL \& Implicit World Modeling with Large Language Models (Show that most LLM-based “world models” are implicit and parameter-encoded, and these paper report limitations that align with our narrative)
%     \item Structured and Symbolic Environment Representations \& Discrete Event Systems Specification (Position DEVS as a mature but under-connected framework in ML/LLM research)
%     \item Automatic Generation and Evaluation of World Models (Justify that our generation and evaluation framework are necessary and novel)
% \end{itemize}
% Looks like there are quite a few papers studying using LLM to generate PDDL domain, and the use it for planning. We should emphasize the difference, i.e., DEVS additionally models time, concurrency, hierarchical, etc.
% }

% World models are commonly used as internal predictive mechanisms that support planning, policy learning, and long-horizon reasoning. Much of the literature represents dynamics implicitly (e.g., learned predictors in latent or observation space) and evaluates models via prediction accuracy or downstream task performance.

% Our setting differs in two respects. First, we constrain attention to \emph{discrete-event} environments where decision-relevant dynamics are best expressed as event ordering and timing, and where an explicit simulator is a natural modeling target. Second, we emphasize \emph{evaluation} via executable behavior and conformance to an operational specification, rather than predictive accuracy against a dataset.

% One line of work adapts LLMs to explicitly predict action outcomes in specific domains. For example, \cite{chae2025wma} introduce world-model–augmented web agents that learn transition-focused abstractions of web environments, allowing agents to simulate the effects of candidate actions before execution and thereby improve robustness in web navigation tasks. A related but more implicit approach treats the LLM itself as both the agent and the world model. In this setting, multi-step reasoning is framed as planning over latent world dynamics encoded in the LLM’s parameters, often augmented with search procedures such as Monte Carlo Tree Search (MCTS) to improve performance on reasoning and planning tasks \cite{hao2023rap}. Extensions of this paradigm further integrate tree search into language agents, assigning the LLM multiple roles—including proposing actions or reasoning steps, providing value estimates, and generating self-reflections—sometimes with external environment feedback to support exploration and correction \cite{pmlr-v235-zhou24r}. Importantly, these approaches do not produce an explicit, executable dynamics model; instead, environment dynamics remain implicit in the LLM’s generative behavior.

% Despite promising empirical results, a growing body of work has identified limitations of implicit world modeling with LLMs. Benchmark evaluations such as PlanBench show that LLMs struggle with reasoning about actions and change in classical planning domains, particularly in multi-step settings \cite{valmeekam2023planbench}. From a conceptual standpoint, \cite{kambhampati2024llmmodulo} argue that LLMs should not be treated as stand-alone planners or verifiers, but rather as approximate knowledge sources integrated with symbolic or otherwise verifiable planning mechanisms. Complementing this view, \cite{katz2024thoughtofsearch} analyze LLM-driven planning and search through the lens of soundness, completeness, and efficiency, noting that unconstrained successor generation can lead to high query costs and weakened guarantees, motivating the incorporation of executable or symbolic constraints to restore reliability.


\subsection{World Models: Explicit and Implicit Approaches}

World models have long served as predictive mechanisms to support planning and reinforcement learning, enabling agents to reason about the consequences of actions before execution and to generate simulated experience for improved learning efficiency \cite{ha2018recurrent}. Broadly, existing work on world modeling can be categorized into two complementary approaches: \emph{explicit world models}, in which environment dynamics are manually implemented inside executable simulators, and \emph{implicit world models}, in which dynamics are captured by learned models, often neural networks.

Several works construct world models as human-designed, executable simulators with well-defined state representations and transition dynamics. For example, \cite{xiang2023language} utilize a robotic simulator as a world model and demonstrate that interaction with such a model can enhance the performance of language models on downstream tasks. Similarly, \cite{feng2025web} present \emph{Web World Models}, a carefully web-based simulator that captures the structure and dynamics of web environments to support controlled agent interaction and evaluation. These approaches emphasize executability, reproducibility, and precise control over environment dynamics.
In parallel, recent work has explored whether environment dynamics can be modeled implicitly by large neural networks, particularly LLMs. In this paradigm, the world model is not represented as a standalone simulator; instead, environment dynamics are encoded in the model’s parameters and accessed through prediction during inference. LLMs are prompted or trained to predict future observations or outcomes conditioned on past events, and these predictions are used to guide action selection or planning \cite{hao2023reasoning,gu2024your}. This approach has been applied across a range of domains, including web environments \cite{gu2024your,feng2025web,chae2025wma}, video games and real-world video data \cite{genie3}, and robotic tasks involving visual, force, or other sensory feedback \cite{wu2023daydreamer}.

These two lines of work reflect different trade-offs in how world dynamics are represented and accessed. Explicit world models provide executable, transparent representations of environment dynamics but typically require substantial manual effort to design and adapt, while implicit neural world models offer greater flexibility and scalability by leveraging learned representations, at the cost of reduced explicitness and multi-turn consistency.



% A complementary line of work focuses on generating structured representations of environment dynamics from natural language, such as planning domains, system models, or simulation artifacts. These representations are typically easier to inspect and can admit rule-based checking.

% This paper aligns with this direction but targets a specific class of structured models: \emph{DEVS discrete-event simulators} that expose an explicit interface for interventions and emit standardized event traces. This interface enables black-box evaluation by executing the simulator under controlled scenarios.

% \textbf{TODO (citations):} Include symbolic world model generation benchmarks and NL-to-system-model work such as \cite{hu2025text2world,jin2025system,article}.
\subsection{Symbolic and Programmatic World Model Generation}

% Another closely related line of work focuses on generating structured representations of environment (transition) dynamics (e.g., planning domains or system models) from natural language to improve interpretability and enable rule-based or formal checking. Regarding symbolic “world/domain models,” existing methods map natural language to executable planning representations (e.g., PDDL), using execution behavior and/or outcomes as key evaluation signals \cite{hu2025text2world,zhang2024open_domain_planning_repr,oswald2024llm_planning_domain_generators}. Meanwhile, other efforts leverage environment interaction to reduce reliance on expert correction, enabling more automated NL to PDDL conversion \cite{mahdavi2024pddl_interaction}.

A closely related line of work focuses on generating structured representations of environment dynamics from natural language, with the goal of improving interpretability and enabling rule-based or formal verification. In the context of symbolic world or domain models, several approaches translate natural-language descriptions into executable planning representations such as PDDL, using plan execution behavior or task outcomes as evaluation signals \cite{hu2025text2world,zhang2024open_domain_planning_repr,oswald2024llm_planning_domain_generators}. Complementary efforts reduce reliance on expert supervision by incorporating environment interaction to iteratively refine generated planning domains \cite{mahdavi2024pddl_interaction}.


% Beyond planning domains, related efforts extend to broader systems-engineering model generation, e.g., generating SysML models from natural-language requirements and evaluating them systematically \cite{jin2025sysmbench,acm2023_nl2sysml_requirements,acm2025_llm_sysml_behavior}. Overall, these studies reflect a trend toward generating explicit world models, while also highlighting the need for more systematic, execution-centered and unified evaluation protocols to mitigate issues such as evaluation randomness and metric indirectness \cite{hu2025text2world,oswald2024llm_planning_domain_generators}.

Beyond planning-oriented representations, related work has explored generating system-level models from natural-language requirements in software and systems engineering. In particular, several studies synthesize SysML artifacts, such as structural diagrams or behavioral models, and evaluate them using consistency checks, requirement coverage, or simulation-based criteria \cite{jin2025sysmbench,acm2023_nl2sysml_requirements,acm2025_llm_sysml_behavior}. These approaches emphasize requirements traceability, design validation, and alignment with established engineering workflows, but typically focus on static structure or high-level behavioral descriptions rather than executable environment dynamics.


Our work aligns with these efforts in seeking explicit, programmatic representations of environment dynamics from natural language, but differs in its choice of modeling formalism and evaluation focus. Compared to PDDL-based planning models, which emphasize symbolic state transitions and goal satisfaction in discrete steps, DEVS provides an execution-oriented semantics centered on discrete events, explicit time, and concurrent interactions, making it well suited for domains such as queues, workflows, distributed systems, and message-driven multi-agent settings. Relative to SysML-based approaches, which primarily support high-level system specification and design reasoning, DEVS offers precise simulation semantics and inherent executability, producing structured event traces that enable black-box, trace-based evaluation under controlled scenarios. Together, these properties position DEVS as a natural foundation for synthesizing and evaluating discrete-event world models from natural-language descriptions.

\subsection{LLM-Based Code Generation and Long-Horizon Development}
% Recent work studies the ability of LLMs and coding agents to implement software from natural language requirements, fix issues in existing repositories, or generate multi-file projects. Evaluation often focuses on functional correctness with respect to unit tests, issue resolution outcomes, or repository-level benchmarks.

% Our problem shares the need for long-horizon code generation, but differs in the evaluation target: we require not only that code executes, but that it conforms to a time-sensitive event semantics specified by an operational contract (CLI parameters, stdin interventions, and a JSONL event schema). This motivates an evaluation approach closer to conformance testing than to generic software unit testing.

% \textbf{TODO (citations):} Add coding-agent evaluation references such as \cite{jimenez2024swebenchlanguagemodelsresolve,ding2025nl2repobenchlonghorizonrepositorygeneration,li2024promptinglargelanguagemodels}.

Recent research has extensively explored the capabilities of large language models and code agents for natural language driven software implementation, fixing issues in existing code repositories, and generating multi-file projects. Evaluations typically focus on unit-test-based functional correctness, issue-fixing success rates, or repository-level benchmarks. Examples include unit-test-driven functional correctness evaluation \cite{chen2021evaluating_code_llms}, repository-repair benchmarks based on real GitHub issues \cite{jimenez2024swebench}, and repository-level multi-file evaluation settings \cite{liu2024repobench,ding2025nl2repo_bench}.

The problem addressed in this paper likewise involves long-horizon code development: coordinating across multiple steps and files, and leveraging execution and tool feedback to iteratively advance the implementation. Existing work often summarizes complex processes with automatable, reproducible outcome metrics (e.g., test pass/fail or task completion rate) and evaluates systems on multi-step interactive software-engineering tasks or repository-level feature-implementation settings \cite{yang2024sweagent,ding2025nl2repo_bench,li2025fea_bench}.

Our evaluation objectives differ fundamentally from these approaches. We require not only executable code but also conformance to a time-sensitive event-semantics specification, defined via standard-input interventions and a unified JSONL event format. Accordingly, our evaluation is closer to conformance testing, where correctness is judged by whether observable input/output behavior satisfies the specification, rather than by generic unit tests \cite{tretmans2006ioco}. It also aligns with the runtime-verification perspective of checking specifications against execution traces \cite{leucker2009rv}.

\subsection{Verification, Validation, and Conformance for Discrete-Event Models}
% Verification and validation of discrete-event and DEVS models has a long history, including static checks, simulation-based validation, and formal analysis of model composition. These methods provide principled ways to detect specification violations.

% We adopt the perspective that execution traces can support systematic evaluation, but we focus on a distinct setting: the DEVS models are synthesized automatically by an LLM-based generator. As a result, the evaluation harness must (i) be robust to a wide range of implementation choices, (ii) provide actionable feedback when specification violations occur, and (iii) operate without requiring code-level equivalence to a single reference implementation.

% \textbf{TODO (citations):} Add DEVS verification/validation references such as \cite{lee2025automated}.


Verification and validation (V\&V) of discrete-event models and DEVS models have a long research history, including static consistency checks, simulation-based verification, and formal analysis of model composition. These approaches provide systematic tools for detecting specification violations. For instance, structural consistency verification can validate the structural definitions of coupled models prior to simulation \cite{lee2025devs_struct_verify}; behavioral-level verification is often instantiated as simulation-driven unit tests or semantic-consistency tests \cite{henares2020unit_test_devs,li2011devs_testing_framework}; and work that integrates DEVS with model checking and applies mixed V\&V in hierarchical/multi-resolution compositional modeling illustrates a path that combines formal analysis with compositional verification \cite{zeigler2017devs_modelchecking,gholami2017constrained_devs_mc}.

In line with this tradition, execution trajectories can serve as a direct basis for systematic evaluation: given input drivers and observation schemes, the runtime traces of the system under test can be compared against specified properties to determine whether the expected behavior is satisfied. This perspective is commonly framed as runtime verification (RV) in the broader software and systems verification literature, where specifications are checked online or offline against event/state sequences \cite{leucker2009rv}. In discrete-event simulation, similar trace-based checking can be achieved by explicitly modeling properties and performing on-demand checks during simulation \cite{dasilva2011simulation_purposes}.

This paper similarly adopts trace-based system evaluation, but under a distinct context: the DEVS models are automatically synthesized by LLM-driven generators, resulting in highly diverse implementation forms. Consequently, the evaluation framework must center on interfaces and observable behaviors to accommodate diverse implementations, and provide actionable localization signals when specification violations are detected. Furthermore, evaluation should target property/behavioral conformance rather than relying on code-level equivalence to a single reference implementation. These requirements naturally connect to black-box testing frameworks and their diagnostic capabilities for failing cases \cite{mclaughlin2020devs_scripting}. They also echo metamorphic testing, which mitigates the oracle problem via metamorphic relations when reliable test oracles are unavailable \cite{chen1998metamorphic}.